% !PRINTED 403
\weberpage{403}

Accordingly, from (2) there results:

\weq{5}{\begin{aligned}
\mathrm{Lim}_{t=\infty}\frac{Z_t}{t} &= \frac{\log\varepsilon}{a\sqrt D}, & D&>0,\\
&= \frac{2\pi}{a\sqrt{-D}}, & D&<0.
\end{aligned}}
But according to Bd. II, § 196, 4. we have

\[ \mathrm{Lim}_{s=1}\sum\frac{s-1}{(Ny)^s} = \mathrm{Lim}_{t=\infty}\frac{Z_t}{t}, \]
and thus from (5) there results:

\weq{6}{\begin{aligned}
\mathrm{Lim}_{s=1}\sum\frac{s-1}{(Ny)^s} &= \frac{1}{a\sqrt D}\log\varepsilon, & D&>0,\\
&= \frac{2\pi}{a\sqrt{-D}}, & D&<0.
\end{aligned}}
Here $y$ runs through the sequence of the integers of the order
$[Q]$ that are divisible by a primary ideal $\mathfrak{a}$ coprime to $Q$,
where, however, of a family of associated numbers only a single
representative, which may now be chosen arbitrarily, is to be
retained (for negative $D$ two or four or six).

Here, associated numbers (with respect to the order $[Q]$) means numbers
whose quotient is \gesperrt{a unit of the order $Q$.}

We now wish to exclude from the numbers $y$ all those that
are not relatively prime to $Q$. Then the sum (6) is modified in
the following way: If $r$ is a prime factor of $Q$, then all numbers
of the order $[Q]$ that are not coprime to $r$ are divisible by $r$,
since they are indeed congruent to a rational number with respect
to the modulus $r$.

If we therefore set

\weq{7}{y = r\,y_1,}
then $y_1$ is divisible by $\mathfrak{a}$, but it need not belong to the order $[Q]$;
it need only belong to the order $\left[\dfrac{Q}{r}\right] = [Q_1]$. If, for
positive discriminant, $\varepsilon_1$ is the fundamental unit of the order
$[Q_1]$, then

\weq{8}{\varepsilon = \varepsilon_1^{\lambda}}
is a power of $\varepsilon_1$, and we have $\sqrt D = r\sqrt{D_1}$. The numbers

\weq{9}{y = r\,y_1\,\varepsilon_1^{t}, \quad t = 0, 1, 2, \ldots, \lambda - 1}
%CONT>>
